\chapter{Probability and Conditional Probability}

A clinic advertises a screening test for a rare disease: only one person in a thousand in
that region has it. The test is good. Out of every hundred sick people, ninety-nine come
back positive. Out of every hundred healthy people, only one is wrongly accused of being
positive. The advertisement writes both numbers as ``99 percent accuracy''.

Now the positive result is in our hands. What is the probability that we are really sick?
In a population of $100{,}000$ people, about $100$ have the disease and $99$ of them are
positive. The rest, $99{,}900$ healthy people, contribute $999$ false positives. So there
are $1.098$ positive results and only $99$ that are right: $99/1.098=11/122$,
that is $0.0902$. About one in eleven, not ninety-nine in a hundred.

Both numbers are right and both answer different questions. The advertisement gives the
probability of a positive result if the person is sick. The patient wants to know the
probability that the person is sick if the result is positive. Swapping those two
directions is not an arithmetic mistake; it is a mistake in naming the quantity, and the
consequence is often a factor of ten.

Chapter 2 already gave us a language for events: union, intersection, complement. Chapter 5
gave us a way to sum infinitely many terms. Neither is enough, because a set carries no
weight: $A\cup B$ is the same set whether $A$ is almost certain or almost certainly not.
What we need is three things: a weight function with rules that must not be broken, a way
to count weights, and a rule for updating weights when new information arrives. This
chapter builds the three, then checks them with a seeded simulation.

\begin{untukanak}
Imagine the guard at a school gate checking bags with a device that beeps. The device is
good: out of a hundred bags that really carry a forbidden object, ninety-nine make it
beep. But the device also beeps for one out of a hundred bags whose contents are perfectly
ordinary. Now suppose that at that school only one bag in a thousand really carries a
forbidden object.

The question is simple and surprising. If the device beeps, is that bag certainly
dangerous? It turns out not. Dangerous bags are indeed rare, while ordinary bags are very
many, so the false beeps from ordinary bags far outnumber the true ones. A guard who looks
only at the red light will hold back many innocent children.

The question that makes us curious: what fraction of all those beeps are right? The answer
can be computed without guessing, as long as we know how many bags there are of each kind.
In this chapter we shall learn why a device that is ``ninety-nine percent accurate'' can
turn into a false accusation in nine cases out of ten, and why the direction of the
question determines the answer.
\end{untukanak}

\section{Probability axioms and their consequences}
\tingkat{basic}

Before a single number is computed, we have to agree on what object is being weighted. In
everyday language: probability is the weight we place on each possibility, under the rule
that the total weight of all possibilities is one and no weight is ever negative. That rule
does not explain the philosophical meaning of probability; it fixes the arithmetic that
every computation has to obey.

\kamus{probability}{peluang}{The weight placed on each possibility, with the total weight
of all possibilities equal to one}

\index{probability}
\index{sample space}

\begin{definisi}[Probability space]
Let $\Omega$ be the set of outcomes, $\mathcal{F}$ a collection of events, that is, the
subsets of $\Omega$ that contain $\Omega$ and are closed under complement and countable
union, and $P\colon\mathcal{F}\to\mathbb{R}$. The function $P$ is called a
\emph{probability} if
\begin{enumerate}
  \item $P(A)\ge 0$ for every $A\in\mathcal{F}$;
  \item $P(\Omega)=1$;
  \item if $A_{1},A_{2},\dots$ are pairwise disjoint, then
        $P\bigl(\bigcup_{i\ge 1}A_{i}\bigr)=\sum_{i\ge 1}P(A_{i})$.
\end{enumerate}
\end{definisi}

The third axiom, countable additivity, separates modern probability theory from the
gambling computations of the seventeenth century. Without it, a probability series such as
$\sum_{n\ge 1}2^{-n}=1$ has no foundation, even though that series is exactly what appears
when we ask on which toss the first success occurs. This axiom is also what allows
probability to be used in limits, the subject of Chapter 4 and Chapter 5.

Three immediate consequences are used over and over, so they are worth proving once.

\begin{proposisi}[Consequences of the axioms]
\index{probability axioms}
For events $A$ and $B$ in a probability space we have
$P(A^{c})=1-P(A)$, $P(\varnothing)=0$, and if $A\subseteq B$ then
$P(A)\le P(B)$. Moreover
$P(A\cup B)=P(A)+P(B)-P(A\cap B)$.
\end{proposisi}

\begin{proof}
Since $A$ and $A^{c}$ are disjoint and their union is $\Omega$, the second and third
axioms give $1=P(\Omega)=P(A)+P(A^{c})$. For the second claim, take $A=\varnothing$ in the
first result. For the third, write $B$ as the disjoint union of $A$ and
$B\setminus A$; then $P(B)=P(A)+P(B\setminus A)\ge P(A)$ because probability is not
negative. For the last, split $A\cup B$ into three disjoint parts $A\setminus B$,
$A\cap B$, and $B\setminus A$. In the same way $A$ and $B$ each split into two parts, so
that
\[
P(A\cup B)=P(A\setminus B)+P(A\cap B)+P(B\setminus A)
=P(A)+P(B)-P(A\cap B).\qedhere
\]
\end{proof}

That last property, inclusion--exclusion, is the reason ``or'' may not simply be added.
For example, in a class of $30$ students, $18$ like coffee and $12$ like tea, and $8$ like
both. Then the probability that a randomly chosen student likes coffee or tea is
\[
\frac{18+12-8}{30}=\frac{22}{30}=0.7333,
\]
not $30/30=1$ as the naive sum $18+12=30$ would give. Eight people are counted twice, and
inclusion--exclusion subtracts them once.

For three events the formula extends with alternating signs:
\[
P(A\cup B\cup C)=P(A)+P(B)+P(C)-P(A\cap B)-P(A\cap C)-P(B\cap C)+P(A\cap B\cap C).
\]
The proof applies the two-event version twice, and that is the first exercise of this
chapter. The pattern of signs is not decoration: whatever was subtracted too often has to
be added back.

One warning for Chapter 18. In a finite space where every outcome has positive probability,
$P(A)=0$ means that $A$ is empty; in a continuous space it does not, since a single point
can have probability zero without being impossible. This chapter works in finite spaces.

\begin{contoh}
Two dice are thrown. What is the probability that the sum of the faces is $8$ or that both
faces are equal? With $36$ equally likely outcomes, the first event contains $5$ outcomes,
the second contains $6$, and their intersection is one outcome, namely $(4,4)$. So the
probability is $(5+6-1)/36=10/36=0.2778$. Without subtracting the intersection the answer
is $11/36=0.3056$, too large by $1/36$.
\end{contoh}

\section{Counting and equally likely outcomes}
\tingkat{intermediate}

If all outcomes are equally likely, the probability of an event is just the ratio of the
number of favorable outcomes to the number of all outcomes. In everyday language: divide
the number of ways we want by the number of ways there are. That definition is not a law of
nature; it is an assumption about symmetry, and the assumption has to be checked, not
presupposed.

\index{counting}
\index{combination}

\kamus{combination}{kombinasi}{A choice of several objects without regard to their order,
like taking five cards from one deck}

\begin{proposisi}[Product rule and counting]
\index{product rule}
If $|A|=m$ and $|B|=n$, then $|A\times B|=mn$. The number of sequences of length $k$ from
$n$ distinct objects without repetition is $n!/(n-k)!$, and the number of subsets of size
$k$ is
\[
\binom{n}{k}=\frac{n!}{k!\,(n-k)!}.
\]
In particular $\sum_{k=0}^{n}\binom{n}{k}=2^{n}$.
\end{proposisi}

\begin{proof}
For the product, every pair $(a,b)$ is determined by one choice from $m$ and one choice
from $n$. For sequences there are $n$ choices for the first place, $n-1$ for the second,
and so on up to $n-k+1$. For subsets, every set of size $k$ can be ordered in $k!$ ways, so
the number of sequences divided by $k!$ gives the number of sets. The last identity counts
all subsets in two ways: summing by size, or choosing in or out for every object.
\end{proof}

Neat counting often shortens a long computation. A full house in poker, three cards of one
rank plus a pair of another rank, is a good example: choose the rank of the triple ($13$
ways), choose three of its four suits ($\binom{4}{3}=4$), choose the rank of the pair ($12$
ways), choose two of its four suits ($\binom{4}{2}=6$). The total is
$13\cdot 4\cdot 12\cdot 6=3.744$ out of $\binom{52}{5}=2{,}598{,}960$ hands, so the
probability is $0.001441$, that is one hand in $694.2$. In $694$ hands the expected number
of full houses is $0.9998$, so we wait about seven hundred hands to see the first one.

Counting also explains events that feel miraculous. Take the birthday problem: $k$ people
in one room, every birthday equally likely among $365$ days, and birthdays independent
between people. The probability that no two people share a birthday is
\[
P(\text{no pair})=\frac{365\cdot 364\cdots (365-k+1)}{365^{k}},
\]
so the probability that there is a pair is its complement.

\index{birthday paradox}

\begin{center}
\small
\begin{tabular}{rrr}
\hline
$k$ & $P(\text{no pair})$ & $P(\text{a pair})$ \\
\hline
$10$ & $0.8831$ & $0.1169$ \\
$23$ & $0.4927$ & $0.5073$ \\
$41$ & $0.0968$ & $0.9032$ \\
$50$ & $0.0296$ & $0.9704$ \\
$57$ & $0.0099$ & $0.9901$ \\
\hline
\end{tabular}
\end{center}

At $k=23$ the probability passes one half, and at $k=50$ it is already $0.9704$. The source
of the confusion is the wrong unit: what collides is not people but pairs of people, and
$23$ people form $\binom{23}{2}=253$ pairs. Every pair matches with probability $1/365$, so
the expected number of pairs is $253/365=0.69315$. If collisions are treated as rare and
independent, the probability of no collision approaches $e^{-0.69315}=0.5000$, differing by
$0.0073$ from the exact value $0.4927$. That approximation is not a proof; it explains why
the threshold falls near twenty-three.

The counting above assumes every birthday equally likely. In real populations that
inequality actually raises the probability of a collision; the symmetry assumption here is
an approximation that has to be checked.

Counting cannot answer the clinical question from before. The outcomes ``sick'' and
``healthy'' are not equally likely: their ratio is one to nine hundred ninety-nine. What is
needed is not a way to count outcomes but a way to use the information that someone has
already been declared positive. That is conditional probability.

\section{Conditional probability and the product rule}
\tingkat{intermediate}

New information narrows the world. If we know that $B$ occurred, the relevant world is no
longer $\Omega$ but $B$, and all weights have to be read again inside that narrower world.
In everyday language: conditional probability is the probability recomputed after part of
the possibilities has been crossed out.

\kamus{conditional probability}{peluang bersyarat}{The probability recomputed after we know
that part of the possibilities no longer applies}

\index{conditional probability}

\begin{definisi}[Conditional probability]
If $A$ and $B$ are events with $P(B)>0$, the probability of $A$ given $B$ is
\[
P(A\mid B)=\frac{P(A\cap B)}{P(B)}.
\]
\end{definisi}

Dividing by $P(B)$ is the choice that makes conditional probabilities sum to one over $B$.
If we also demand $P(B\mid B)=1$ and that ratios between events inside $B$ do not change,
that division is the only possible choice; hence this definition is not an arbitrary
convention.

The narrowing can be seen directly in a cross table. Figure~\ref{fig:silang} arranges the
clinical population into four cells. Before any result, the denominator is all $100{,}000$
people and $P(S)=100/100{,}000=0.001$. After a positive result is known, only the $+$
column, which holds $1.098$ people, is still counted, and
$P(S\mid+)=99/1.098=0.0902$. The person does not change; only the set that may still be
counted changes.

    \begin{figure}[htbp]
\centering
\begin{tikzpicture}[font=\small, x=1cm, y=1cm]
  \fill[accent!10] (1.95,1.05) rectangle (3.25,-1.1);
  \draw[layergrey, thick] (1.9,1.05)--(7.3,1.05);
  \draw[layergrey] (1.9,0.35)--(7.3,0.35);
  \draw[layergrey] (1.9,-0.35)--(7.3,-0.35);
  \draw[layergrey, thick] (1.9,-1.1)--(7.3,-1.1);
  \node[anchor=east] at (1.75,1.4) {test result};
  \node at (2.6,1.4) {$+$};
  \node at (4.3,1.4) {$-$};
  \node at (6.5,1.4) {total};
  \node[anchor=east] at (1.75,0.7) {sick $S$};
  \node at (2.6,0.7) {99};
  \node at (4.3,0.7) {1};
  \node at (6.5,0.7) {100};
  \node[anchor=east] at (1.75,0) {healthy $S^{c}$};
  \node at (2.6,0) {999};
  \node at (4.3,0) {98{,}901};
  \node at (6.5,0) {99{,}900};
  \node[anchor=east] at (1.75,-0.7) {total};
  \node at (2.6,-0.7) {1.098};
  \node at (4.3,-0.7) {98{,}902};
  \node at (6.5,-0.7) {100{,}000};
  \draw[->, accent, thick] (2.6,-1.65)--(2.6,-1.15);
  \node[anchor=north, align=center, text width=10.4cm, font=\footnotesize] at (4.6,-1.95)
    {once $+$ is known, only this column is counted:
     $99/1.098=11/122=0.0902$};
\end{tikzpicture}
\caption{Cross table of a screening test on $100{,}000$ people. Rows are the true state,
columns are the test result. Once a positive result is known, the sample space narrows
from the whole table to the $+$ column, which holds $1.098$ people, and $99$ of them
are really sick.}
\label{fig:silang}
\end{figure}    

\begin{proposisi}[Conditional probability is a probability]
\index{conditional probability}
For fixed $B$ with $P(B)>0$, the function $A\mapsto P(A\mid B)$ satisfies the three axioms
of probability on $\mathcal{F}$.
\end{proposisi}

\begin{proof}
Nonnegativity comes from $P(A\cap B)\ge 0$. Normality comes from $P(\Omega\mid B)=1$,
because $P(B)/P(B)=1$. For additivity, if $A_{1},A_{2},\dots$ are pairwise disjoint then
$A_{i}\cap B$ are also pairwise disjoint, so that
\[
P\Bigl(\bigcup_{i}A_{i}\mid B\Bigr)
=\frac{P\bigl(\bigcup_{i}(A_{i}\cap B)\bigr)}{P(B)}
=\frac{\sum_{i}P(A_{i}\cap B)}{P(B)}
=\sum_{i}P(A_{i}\mid B).\qedhere
\]
\end{proof}

Hence all consequences of the axioms hold for conditional probability: complement,
monotonicity, and inclusion--exclusion remain valid when every $P$ is replaced by
$P(\cdot\mid B)$. What does not hold is swapping the direction of the condition.

From the definition the product rule is born at once
\[
P(A\cap B)=P(A\mid B)\,P(B)=P(B\mid A)\,P(A),
\]
which extends to
$P(A_{1}\cap\cdots\cap A_{n})=P(A_{1})P(A_{2}\mid A_{1})\cdots
P(A_{n}\mid A_{1}\cap\cdots\cap A_{n-1})$. This rule reads a staged experiment as a
product of probabilities at each stage.

A single deck of $52$ cards gives an example of a direction that is easy to swap. Let $A$
be the event that the first card is a heart and $B$ the event that the second card is red.
Then
\[
P(A\cap B)=\frac{13\cdot 25}{52\cdot 51}=\frac{25}{204},
\]
so that
\[
P(A\mid B)=\frac{25/204}{1/2}=\frac{25}{102}=0.2451,
\qquad
P(B\mid A)=\frac{25/204}{1/4}=\frac{25}{51}=0.4902 .
\]
Two numbers that differ by a factor of two, from the same data. The mistake is not in the
arithmetic but in the placement of the events to the left and right of the vertical bar.

\begin{bacakeras}
Two voices. The first does not know yet, the second already knows. Read them in turn.

First voice: ``My world is as wide as a playing field. Every door is possible. I guess with
my eyes closed.''

Second voice: ``My world is one room. One door is already open, and I count only what is
inside it.''

First voice: ``If I know something, does my probability change?''

Second voice: ``It is not your probability that changes. It is the world you count that
changes. The world shrinks after we learn something: weight that was spread out now
gathers in the places that are still possible. The denominator narrows, and the number that
comes out sounds new even though the arithmetic has not changed since the beginning.''

First voice: ``So it is the denominator that moves?''

Second voice: ``The denominator moves, and what is left to be recounted is the remainder
that may still take part.''

Repeat together, three times, more and more slowly: ``The world shrinks, the probability
does not betray us. Only what may still be counted changes.''
\end{bacakeras}

The second tool splits the population into several groups that do not overlap. If
$B_{1},\dots,B_{k}$ are pairwise disjoint and their union is $\Omega$, then
\[
P(A)=\sum_{i=1}^{k}P(A\mid B_{i})\,P(B_{i}).
\]

\begin{proposisi}[Law of total probability]
\index{total probability}
If $B_{1},\dots,B_{k}$ are pairwise disjoint, their union is $\Omega$, and every
$P(B_{i})>0$, then $P(A)=\sum_{i}P(A\mid B_{i})P(B_{i})$ for every event $A$.
\end{proposisi}

\begin{proof}
The events $A\cap B_{1},\dots,A\cap B_{k}$ are pairwise disjoint and their union is $A$.
So $P(A)=\sum_{i}P(A\cap B_{i})$, and every term is rewritten with the product rule as
$P(A\mid B_{i})P(B_{i})$.
\end{proof}

This law looks like a formality until it is used on the clinical example: it is what turns
two numbers from an advertisement into one useful number, and that is the content of the
next section.

Before that, one example showing that ordinary sentences and mathematical events do not
always correspond. A family has two children, each child a boy or a girl with equal
probability, and the sexes of the two children are independent. If we know that
\emph{at least one} child is a boy, the probability that both are boys is $1/3=0.3333$: of
the four equally likely arrangements three remain, and only one fits. If what we know is
that \emph{the older child} is a boy, the probability is $1/2=0.5$. Two similar sentences
give different answers, because the first excludes one arrangement and the second two.

\section{Bayes' rule and the interpretation of test results}
\tingkat{intermediate}

Bayes' rule reverses the direction of the condition: from the probability of a result given
sickness it computes the probability of sickness given the result. In everyday language: of
all the ways a result can appear, what fraction comes through the path we are asking about.

\index{Bayes' rule}
\index{prevalence}

\begin{teorema}[Bayes' rule]
If $B_{1},\dots,B_{k}$ are pairwise disjoint with union $\Omega$ and $P(B_{i})>0$, then
for every event $A$ with $P(A)>0$,
\[
P(B_{j}\mid A)=\frac{P(A\mid B_{j})\,P(B_{j})}
{\sum_{i=1}^{k}P(A\mid B_{i})\,P(B_{i})}.
\]
\end{teorema}

\begin{proof}
By definition, $P(B_{j}\mid A)=P(A\cap B_{j})/P(A)$. The numerator is
$P(A\mid B_{j})P(B_{j})$ by the product rule, and the denominator is $P(A)$ by the law of
total probability. So the formula contains no new assumption; it is the definition of
conditional probability written twice, in two different directions.
\end{proof}

Hence there is no ``Bayesian method'' standing on top of the axioms; what there is, is
bookkeeping. The numerator is the path we are asking about, the denominator is the sum of
all paths that produce the same result.

\begin{resep}
\begin{enumerate}
  \item First name two things: the hypothesis $H$ we are asking about and the evidence $E$
        that has already been seen. Write the question as $P(H\mid E)$, not $P(E\mid H)$.
  \item Write the numerator as a single path: $P(E\mid H)\,P(H)$. That is the weight of the
        world in which $H$ is true.
  \item Write the denominator as the sum of all paths that produce $E$, including the path
        through $H^{c}$: $P(E\mid H)P(H)+P(E\mid H^{c})P(H^{c})$.
  \item Divide, then check for sense: the result lies between $0$ and $1$, and is smaller
        than $P(E\mid H)$ when $H$ is rare.
\end{enumerate}
A warning about order: numerator and denominator both use $P(E\mid\cdot)$. Writing
$P(\cdot\mid E)$ in the denominator is the mistake that happens most often, and the size of
the number does not tell us that the direction has been swapped.
\end{resep}

The tree in Figure~\ref{fig:pohon} writes out every road to a positive result. Each leaf is
the product of the probabilities on its branch, and the two leaves marked $+$ sum to
$P(+)=0.01098$. The Bayes numerator is just the first leaf, while the denominator is the
sum of those two leaves.

    \begin{figure}[htbp]
\centering
\begin{tikzpicture}[font=\small, every node/.style={inner sep=2pt}]
  \node (r) at (0,0) {start};
  \node (s) at (2.2,1.7) {$S$};
  \node (sc) at (2.2,-1.7) {$S^{c}$};
  \node (sp) at (4.6,2.5) {$+$};
  \node (sm) at (4.6,0.9) {$-$};
  \node (scp) at (4.6,-0.9) {$+$};
  \node (scm) at (4.6,-2.5) {$-$};
  \draw[->] (r)--(s) node[midway, above left=-2pt] {$0.001$};
  \draw[->] (r)--(sc) node[midway, below left=-2pt] {$0.999$};
  \draw[->, accent, thick] (s)--(sp) node[midway, above left=-2pt] {$0.99$};
  \draw[->, layergrey] (s)--(sm) node[midway, below left=-2pt] {$0.01$};
  \draw[->, accent, thick] (sc)--(scp) node[midway, above left=-2pt] {$0.01$};
  \draw[->, layergrey] (sc)--(scm) node[midway, below left=-2pt] {$0.99$};
  \node[anchor=west, font=\footnotesize] at (5.0,2.5) {$0.001\cdot 0.99=0.00099$};
  \node[anchor=west, font=\footnotesize] at (5.0,0.9) {$0.001\cdot 0.01=0.00001$};
  \node[anchor=west, font=\footnotesize] at (5.0,-0.9) {$0.999\cdot 0.01=0.00999$};
  \node[anchor=west, font=\footnotesize] at (5.0,-2.5) {$0.999\cdot 0.99=0.98901$};
  \node[anchor=north, align=center, text width=10.6cm, font=\footnotesize] at (4.4,-3.0)
    {$P(+)=0.00099+0.00999=0.01098$, \quad $P(S\mid+)=0.00099/0.01098=0.0902$};
\end{tikzpicture}
\caption{Two-stage tree of a screening test. The first branch separates sick from healthy,
the second separates positive results from negative ones; the product at each leaf is the
probability of that path, and the two positive leaves sum to $P(+)=0.01098$.}
\label{fig:pohon}
\end{figure}    

Back to the clinic. Write $S$ for the event of being sick and $+$ for a positive result.
From the advertisement and the prevalence of the disease we have $P(S)=0.001$,
$P(+\mid S)=0.99$, and $P(+\mid S^{c})=0.01$. The law of total probability gives
\[
P(+)=0.99\cdot 0.001+0.01\cdot 0.999=0.00099+0.00999=0.01098,
\]
so that
\[
P(S\mid +)=\frac{0.00099}{0.01098}=\frac{11}{122}=0.0902 .
\]
The false positive term, $0.00999$, is ten times larger than the true detection term,
$0.00099$; their ratio is $999/99=10.09$. Not because the test is bad: healthy people are
$999$ times more numerous, $99{,}900$ against $100$.

A change of prevalence changes the answer far more strongly than a change in the quality of
the test. With sensitivity and specificity fixed at $0.99$, the following table shows
$P(S\mid +)$ for several prevalences.

\begin{center}
\small
\begin{tabular}{rrrr}
\hline
$P(S)$ & $P(S\mid +)$ & $1/P(S\mid +)$ & reading \\
\hline
$0.0001$ & $0.0098$ & $102.0$ & one in a hundred and two \\
$0.001$ & $0.0902$ & $11.09$ & one in eleven \\
$0.010$ & $0.5000$ & $2.00$ & one in two \\
$0.100$ & $0.9167$ & $1.09$ & eleven in twelve \\
\hline
\end{tabular}
\end{center}

At a prevalence of $1$ percent the same test gives $P(S\mid+)=1/2$ exactly, not $0.99$.
At a prevalence of $10$ percent, $P(S\mid+)=11/12=0.9167$. These numbers explain a health
policy that is often criticized: mass screening for a rare disease produces many false
positives, and for that reason a positive result is usually followed by a second test
rather than by treatment.

The second test uses the same rule, with the prior already updated. If the second test is
independent of the first and of the same quality, then
\[
P(S\mid +,+)=\frac{0.99\cdot (11/122)}{0.99\cdot (11/122)+0.01\cdot (111/122)}
=\frac{10.89}{10.89+1.11}=\frac{363}{400}=0.9075 .
\]
Note that $0.9075$ is neither $0.99$ nor $0.9999$. Two positive tests raise the probability
from $0.0902$ to $0.9075$, almost ten times, but they do not yet give certainty. What
changes is only the prevalence inside the group that is already positive; the quality of
the test is unchanged. With this quality of $0.99$, the posterior exceeds $0.99$ only when
the initial prevalence is already past $1/2$.

One technical point: $P(S\mid+)=0.0902$ applies to a population with prevalence $1/1000$,
not to a patient who arrives with symptoms. Symptoms change the prevalence first, and that
is why physicians ask about history before interpreting a laboratory result.

The last example of this section releases Bayes from the clinical test. Three bags hold ten
marbles each: bag $A$ one red, $B$ five red, $C$ nine red. We choose a bag at random, draw
one marble, and the one that comes out is red. Before the draw the three bags each have
probability $1/3$; afterwards they have $0.0667$, $0.3333$, and $0.6$, as in
Figure~\ref{fig:bayes-batang}. No weight is lost; weight merely moves to the bag that more
easily produces that evidence.

    \begin{figure}[htbp]
\centering
\begin{tikzpicture}[font=\small, x=1cm, y=1cm]
  \draw[->, layergrey] (0,0)--(7.4,0);
  \draw[->, layergrey] (0,0)--(0,4.3);
  \foreach \y/\lab in {0/0, 1.2/0.2, 2.4/0.4, 3.6/0.6} {
    \draw[layergrey] (0,\y)--(-0.08,\y);
    \node[anchor=east, font=\footnotesize] at (-0.12,\y) {\lab};
  }
  \fill[layerindigo] (0.65,0) rectangle (1.20,2.0);
  \fill[layerteal] (1.40,0) rectangle (1.95,2.0);
  \fill[accent] (2.15,0) rectangle (2.70,2.0);
  \fill[layerindigo] (4.35,0) rectangle (4.90,0.4);
  \fill[layerteal] (5.10,0) rectangle (5.65,2.0);
  \fill[accent] (5.85,0) rectangle (6.40,3.6);
  \draw[dashed, layergrey] (0,2.0)--(6.9,2.0);
  \node[font=\footnotesize, anchor=south] at (0.92,2.02) {$1/3$};
  \node[font=\footnotesize, anchor=south] at (1.67,2.02) {$1/3$};
  \node[font=\footnotesize, anchor=south] at (2.42,2.02) {$1/3$};
  \node[font=\footnotesize, anchor=south] at (4.62,0.44) {$0.0667$};
  \node[font=\footnotesize, anchor=south] at (5.37,2.02) {$0.3333$};
  \node[font=\footnotesize, anchor=south] at (6.12,3.62) {$0.6$};
  \node[font=\footnotesize, anchor=north] at (0.92,-0.08) {$A$};
  \node[font=\footnotesize, anchor=north] at (1.67,-0.08) {$B$};
  \node[font=\footnotesize, anchor=north] at (2.42,-0.08) {$C$};
  \node[font=\footnotesize, anchor=north] at (4.62,-0.08) {$A$};
  \node[font=\footnotesize, anchor=north] at (5.37,-0.08) {$B$};
  \node[font=\footnotesize, anchor=north] at (6.12,-0.08) {$C$};
  \node[font=\footnotesize] at (1.67,-0.85) {before};
  \node[font=\footnotesize] at (5.37,-0.85) {after};
\end{tikzpicture}
\caption{Bayes update on three bags. Before any evidence the three bags share the weight
$1/3$; after one red marble is seen, the weights move to $0.0667$, $0.3333$, and $0.6$. The
bag that most easily produces that evidence rises the highest.}
\label{fig:bayes-batang}
\end{figure}    

\section{Independence, and why disjoint is not independent}
\tingkat{advanced}

Two events are called disjoint if they cannot occur together. Two events are called
independent if knowing one does not change the probability of the other. The two sentences
sound similar, and swapping them is the most expensive mistake in this chapter. In everyday
language: disjoint means never together; independent means no news from one to the other.

\kamus{independent}{independen}{Events that send each other no news: knowing one does not
change the probability of the other}

\begin{untukanak}
There are two things that are often swapped, and the difference can be seen without numbers.
At one school, two students are never in class at the same time: one studies in the
morning, one studies in the afternoon. If we know that Andi has come in, we know that Budi
has not. They are never together, and precisely for that reason news about one gives news
about the other.

Now two other students. Both attend the same class, but they sit far apart and work on
different exams. Their grades send each other no news: knowing Sari's grade does not help
us guess Tono's.

The first group is called never together. The second group is called no news to each other.
Both sound like ``no relation'', even though the first is in fact the strongest relation:
news about one determines the other. In this chapter we shall learn why two things that are
never together cannot possibly send each other news, and why swapping the two makes the
computation miss by more than a little.
\end{untukanak}

\index{independence}
\index{disjoint events}

\begin{definisi}[Independence of two events]
Events $A$ and $B$ are \emph{independent} if
\[
P(A\cap B)=P(A)\,P(B).
\]
\end{definisi}

This definition is symmetric, and when $P(B)>0$ it is equivalent to $P(A\mid B)=P(A)$. Note
that the definition does not mention cause and effect: two events can be independent
without one causing the other, and can be dependent without one causing the other.

\begin{proposisi}[Disjoint is not independent]
\index{disjoint events}
If $A$ and $B$ are disjoint and $P(A)>0$ and $P(B)>0$, then $A$ and $B$ are not
independent. More precisely $P(A\cap B)=0<P(A)P(B)$, and $P(A\mid B)=0<P(A)$.
\end{proposisi}

\begin{proof}
Disjoint means $A\cap B=\varnothing$, so $P(A\cap B)=P(\varnothing)=0$. Since
$P(A)P(B)>0$, the equality in the definition of independence fails. For the conditional
statement, $P(A\mid B)=0/P(B)=0$, which is smaller than $P(A)>0$.
\end{proof}

The mistake moves in two directions, and both directions can be computed.

Take one die and two pairs of events with the same marginals but a different relation. First
pair: $A=\{1,2\}$ and $B=\{3,4\}$, disjoint, with $P(A)=P(B)=1/3$. Second pair:
$C=\{1,2,3\}$ and $D=\{1,4\}$, with $P(C)=1/2$ and $P(D)=1/3$. For the first pair,
\[
P(A)P(B)=\frac{1}{9}=0.1111,\qquad P(A\cap B)=0,\qquad P(A\mid B)=0,
\]
so the gap in the definition of independence is $0.1111$ and the gap in conditional
probability is $0.3333$, because $P(A)=1/3$. For the second pair,
\[
P(C)P(D)=\frac{1}{6}=0.1667=P(C\cap D),\qquad P(C\mid D)=\frac{1}{2}=P(C),
\]
so both gaps are zero.

What matters is not the number $1/6$ but the fact that two pairs can have the same marginals
with a completely opposite relation. The pair $E=\{1,2,3\}$ and $F=\{4,6\}$ also has
marginals $1/2$ and $1/3$ but is disjoint. Marginals do not determine the relation; only the
intersection determines it, so every claim of independence is an additional claim about the
world.

For three events, pairwise independence is not enough. The events $A,B,C$ are called
mutually independent if $P(A\cap B)=P(A)P(B)$, $P(A\cap C)=P(A)P(C)$,
$P(B\cap C)=P(B)P(C)$, and $P(A\cap B\cap C)=P(A)P(B)P(C)$.

\begin{contoh}[Pairwise but not mutually independent]
\index{pairwise independence}
Two coins are tossed. Let $A$ be heads on the first coin, $B$ heads on the second coin, and
$C$ both coins the same; each has probability $1/2$. Every
pair intersects in one outcome, so $P(A\cap B)=P(A\cap C)=P(B\cap C)=1/4$, which equals
$1/2\cdot 1/2$. But $A\cap B\cap C$ occurs only when both are heads,
so its probability is $1/4$, whereas $P(A)P(B)P(C)=1/8$. So the three are pairwise
independent and not mutually independent; the gap is $0.125$.
\end{contoh}

A final warning: independence does not survive conditioning. If $A$ and $B$ are independent
with probability $1/2$, then given $A\cup B$, $P(A\mid A\cup B)=2/3$
while $P(A\cap B\mid A\cup B)=1/3\neq 4/9$. Combining two independent events and then
conditioning on their union actually creates dependence, so the sentence ``$A$ and $B$ are
independent'' is always relative to the probability space in use.

\begin{tangan}
Materials: twelve cards, six pairs of equal rank (ace through six), each pair one red and
one black.

Step one, make the pairs. Lay the twelve cards face up in two rows; put each pair of equal
rank side by side, red on the left, black on the right. Check with your eyes: rank sends no
news about color, since every rank has a red and a black. Those two properties send each
other no news.

Step two, filter. Shuffle, cover them all, then turn over one card: say the four of hearts.
Give one hint, ``this card's mate is black'', then set aside every red card that is still
face down. Of the eleven face-down cards, only six are black.

What you must see: the denominator changes from eleven to six, while the cards do not change
at all. The probability of guessing its mate rises from one in eleven to one in six. Two
properties that sent each other no news now send each other news, and that is what one hint
does.
\end{tangan}

\section{Two paradoxes that demand a computation}
\tingkat{advanced}

Paradoxes in probability are almost always a collision between intuition and definition, and
the collision is settled only by computing. The two examples below show different things:
the answer depends on the rules of the game, and on the unit that is counted.

\index{Monty Hall}

Monty Hall. Three closed doors; one holds a car, two hold goats. We choose one door. The
host, who knows what is behind all the doors, opens one other door with a goat, then offers
the chance to switch to the door that is still closed. Does switching raise the
probability?

Compute with the law of total probability over the three possible locations of the car,
each with probability $1/3$, and suppose that we first choose door $1$. If the car is
behind door $1$, switching loses. If the car is behind door $2$, the host must open door
$3$, so switching wins. If the car is behind door $3$, he opens door $2$, and switching
wins. So switching wins in two of the three cases:
\[
P(\text{win by switching})=\frac{2}{3},\qquad
P(\text{win by staying})=\frac{1}{3}.
\]
The gap is $1/3$, not zero.

The argument ``two doors left, so $1/2$'' is wrong because it ignores that the host does not
choose a door freely. He is forbidden to open the door with the car and forbidden to open
the door we chose, so his action carries information about the door he does not open.

That information disappears when the rule changes. Suppose the host opens one of the other
two doors at random, without looking at what is behind it, and the game is cancelled if the
car is revealed. Given that the game runs as far as the offer,
\[
P(\text{car behind door }1\mid \text{goat revealed})
=\frac{1/3}{1/3+1/6+1/6}=\frac{1}{2},
\]
because if the car is behind door $2$ or $3$, he has only probability $1/2$ of showing a
goat. Under that rule switching does not help: ``switching has probability $2/3$'' is not a
property of the doors but a property of the protocol.

\index{birthday paradox}

Birthdays. The previous section already computed $0.5073$ for $23$ people and $0.9704$ for
$50$ people. The paradox is not in the numbers but in the unit: we compare it with the
probability that one particular person matches us, whereas what is computed is $253$
pairs at once. The same mistake appears in hash tables and in DNA match tests.

The prosecutor's fallacy is the same directional mistake in a criminal case. Suppose the
probability that a feature appears by coincidence is one in a million, and the feature is
found in a suspect taken from a database of a million people with no other evidence. Then
$P(\text{match}\mid\text{innocent})=10^{-6}$, but
\[
P(\text{guilty}\mid\text{match})
=\frac{10^{-6}\cdot 1}{10^{-6}\cdot 1+(1-10^{-6})\cdot 10^{-6}}=0.5000003,
\]
that is practically one half, not near certainty. The two numbers $10^{-6}$ are different
quantities: the first is the probability of the feature by coincidence, the second is the
prevalence of the suspect inside the database. Swapping them is another face of the same
mistake, this time in a courtroom.

\begin{tantangan}
Three cards on a table. The first card is red on both sides, the second is black on both
sides, the third is red on one side and black on the other. We shuffle the three face down,
take one, and lay it on the table. The visible side is red. What is the probability that
the side underneath is also red?

Almost everyone answers one half: they say only two cards have a red side, so there are two
possibilities. The correct answer is not one half.

Hint 1: do not count cards, count sides. How many red sides are there in total, and how many
of them sit next to a red side?

Hint 2: the visible red side can come from three different red sides, and two of them are on
the red-red card.

If you are stuck, go back to the section on conditional probability and rewrite the sample
space as six sides, not three cards.
\end{tantangan}

\section{Worked computations}
\tingkat{intermediate}

The three computations below are carried out with exact fractions, so that no rounding hides
a step. All of them have been checked with Python, and the numbers that appear are the same
numbers used in the previous sections.

First, two-step Bayes. With $P(S)=1/1000$, $P(+\mid S)=99/100$, and
$P(+\mid S^{c})=1/100$, the law of total probability gives
\[
P(+)=\frac{99}{100}\cdot\frac{1}{1000}+\frac{1}{100}\cdot\frac{999}{1000}
=\frac{99+999}{100{,}000}=\frac{1098}{100{,}000}.
\]
The numerator is two separate groups: $99$ sick people who are detected and $999$ healthy
people who are wrongly accused. Hence
\[
P(S\mid +)=\frac{99/100{,}000}{1098/100{,}000}=\frac{99}{1098}=\frac{11}{122}=0.090164 .
\]
A second test of the same quality uses $11/122$ as the new prevalence and gives $0.9075$,
just as in the computation of the previous section. The denominator is $1200$, not a number
in the millions: the factor $1/122$ appears in both terms and cancels, a sign that what
changed is only the prevalence, not the quality of the test.

Second, disjoint versus independent, with matched marginals.

\begin{center}
\small
\begin{tabular}{lrrrr}
\hline
pair & $P(\text{first})$ & $P(\text{second})$ & $P(\text{intersection})$ & $P(\text{first}\mid\text{second})$ \\
\hline
$C=\{1,2,3\}$, $D=\{1,4\}$ & $0.5$ & $0.3333$ & $0.1667$ & $0.5$ \\
$E=\{1,2,3\}$, $F=\{4,6\}$ & $0.5$ & $0.3333$ & $0$ & $0$ \\
\hline
\end{tabular}
\end{center}

In the first row $P(C)P(D)=1/6=0.1667$ equals the intersection, so the gap is $0$. In the
second row $P(E)P(F)=1/6$ while the intersection is $0$, so the gap is $0.1667$; the
conditional probability is $0$ while $P(E)=0.5$. Anyone who takes ``cannot happen together''
to mean ``does not influence each other'' will report $0.1667$ where $0$ belongs.

Third, Monty Hall with the protocol stated. If the host always opens a goat door, the
probability of winning by switching is $2/3=0.6667$ and by staying $1/3=0.3333$. If he opens
at random and the game is cancelled when the car is revealed, the game runs with probability
$2/3$, and given that, both strategies are worth $1/2=0.5$. Two protocols, one set of doors,
two answers.

\kertas{For a test with prevalence $1/2000$, sensitivity $0.95$, and specificity
$0.90$, compute $P(S\mid+)$ with the $100{,}000$ table as above, then compute it again with
specificity $0.99$. Write both answers as fractions and as decimals with four digits, and
explain which part changes the most: the numerator or the denominator. After that compute
$P(S\mid +,+)$ for a second test of the same quality, and check whether the result is still
below $0.95$.}

\section{Computational lab: Monte Carlo as a verification tool}
\tingkat{advanced}

All the numbers in this chapter are derived, not measured. Simulation does not replace a
derivation, but it catches the mistakes that happen most often: a swapped direction of the
condition, a wrong denominator, and an independence assumption used without grounds.
Generate the same experiment many times, compute the ratio of successes, then compare it
with the exact value. By the law of large numbers (Chapter 19), that ratio tends to the true
probability; what remains is only how fast, measured by the standard error
$\sqrt{p(1-p)/m}$ with $m$ the number of trials that produce the event.

\index{Monte Carlo}
\index{standard error}

A fixed seed is used so that the output can be checked again word for word. A fixed seed
does not make the result more correct; it only makes it repeatable.

\begin{lab}
\begin{verbatim}
import numpy as np

rng = np.random.default_rng(20261024)

print("== 1. diagnostic test: prevalence 0.1%, sens 99%, spec 99% ==")
n = 1_000_000
sick = rng.random(n) < 0.001
positive = np.where(sick, rng.random(n) < 0.99, rng.random(n) < 0.01)
m = int(positive.sum())
ppv = sick[positive].mean()
err = np.sqrt(ppv * (1 - ppv) / m)
exact = 0.99 * 0.001 / (0.99 * 0.001 + 0.01 * 0.999)
print("sick=%d  positive=%d  true_positive=%d"
      % (sick.sum(), m, int((sick & positive).sum())))
print("PPV empirical=%.6f  exact=%.6f  std error=%.6f  z=%.3f"
      % (ppv, exact, err, (ppv - exact) / err))

print("== 2. Monty Hall, 10^5 games ==")
N = 100_000
car = rng.integers(0, 3, N)
pick = rng.integers(0, 3, N)
opened = np.where(pick == car,
                  (pick + 1 + rng.integers(0, 2, N)) % 3,
                  3 - pick - car)
switch = 3 - pick - opened
win_stay = float((pick == car).mean())
win_switch = float((switch == car).mean())
se_switch = np.sqrt(win_switch * (1 - win_switch) / N)
print("win without switch=%.4f  win switch=%.4f  std error=%.5f"
      % (win_stay, win_switch, se_switch))
print("exact 1/3=%.6f  2/3=%.6f" % (1/3, 2/3))

print("== 3. host opens a door at random ==")
open_random = (pick + 1 + rng.integers(0, 2, N)) % 3
valid = open_random != car
switch_random = 3 - pick - open_random
print("valid games=%d of %d (%.4f)" % (valid.sum(), N, valid.mean()))
print("win switch (valid)=%.4f  without switch (valid)=%.4f"
      % (float((switch_random[valid] == car[valid]).mean()),
         float((pick[valid] == car[valid]).mean())))

print("== 4. birthday probability, 20,000 rooms ==")
for k in (23, 50):
    days = rng.integers(0, 365, size=(20000, k))
    days.sort(axis=1)
    dup = (np.diff(days, axis=1) == 0).any(axis=1)
    q = 1.0
    for i in range(k):
        q *= (365 - i) / 365
    f = float(dup.mean())
    se = np.sqrt(f * (1 - f) / 20000)
    print("k=%2d  empirical=%.4f  exact=%.4f  std error=%.5f  z=%+.2f"
          % (k, f, 1 - q, se, (f - (1 - q)) / se))

print("== 5. disjoint versus independent, 10^6 die throws ==")
N = 1_000_000
die = rng.integers(1, 7, N)
A = np.isin(die, [1, 2])
B = np.isin(die, [3, 4])
C = np.isin(die, [1, 2, 3])
D = np.isin(die, [1, 4])
print("disjoint:  P(A)=%.4f P(B)=%.4f P(AnB)=%.4f P(A|B)=%.4f gap=%.4f"
      % (A.mean(), B.mean(), (A & B).mean(), (A & B).sum() / B.sum(),
         A.mean() * B.mean() - (A & B).sum() / B.sum()))
print("independent:  P(C)=%.4f P(D)=%.4f P(CnD)=%.4f P(C|D)=%.4f gap=%.4f"
      % (C.mean(), D.mean(), (C & D).mean(), (C & D).sum() / D.sum(),
         abs(C.mean() - (C & D).sum() / D.sum())))
print("exact disjoint: 1/3, 1/3, 0, 0, 1/9=%.6f" % (1 / 9))
print("exact independent: 1/2, 1/3, 1/6, 1/2, 0")
\end{verbatim}    
\end{lab}

The actual output, exactly as it came out:

\begin{verbatim}
== 1. diagnostic test: prevalence 0.1%, sens 99%, spec 99% ==
sick=963  positive=11094  true_positive=952
PPV empirical=0.085812  exact=0.090164  std error=0.002659  z=-1.637
== 2. Monty Hall, 10^5 games ==
win without switch=0.3325  win switch=0.6675  std error=0.00149
exact 1/3=0.333333  2/3=0.666667
== 3. host opens a door at random ==
valid games=66472 of 100000 (0.6647)
win switch (valid)=0.4998  without switch (valid)=0.5002
== 4. birthday probability, 20,000 rooms ==
k=23  empirical=0.5019  exact=0.5073  std error=0.00354  z=-1.54
k=50  empirical=0.9701  exact=0.9704  std error=0.00120  z=-0.19
== 5. disjoint versus independent, 10^6 die throws ==
disjoint:  P(A)=0.3333 P(B)=0.3335 P(AnB)=0.0000 P(A|B)=0.0000 gap=0.1112
independent:  P(C)=0.5001 P(D)=0.3334 P(CnD)=0.1667 P(C|D)=0.5000 gap=0.0001
exact disjoint: 1/3, 1/3, 0, 0, 1/9=0.111111
exact independent: 1/2, 1/3, 1/6, 1/2, 0
\end{verbatim}    

Four things deserve reading from that output, and two of them actually show the limits of
simulation.

First, the first line of part 1 reports $963$ sick people out of a million, whereas the
expectation is $1{,}000$ with standard deviation $\sqrt{10^{6}(0.001)(0.999)}=31.6$. So
$963$ lies $1.17$ standard deviations below the expectation: a normal swing, not a
programming oversight. The expected number of positive results becomes
$963\cdot 0.99+999{,}037\cdot 0.01=10{,}943.7$,
and the observed $11{,}094$ is $1.4$ standard deviations away. The empirical conditional
probability $0.0858$ differs from the exact $0.0902$ by $0.0044$, that is $1.64$ standard
errors. That gap is far smaller than the distance from $0.09$ to $1$, so the interpretation
does not change.

Second, Monty Hall gives $0.6675$ and $0.3325$ with standard error $0.0015$, while the
expected gap between the strategies is $1/3$. Agreement to three digits is not luck; the
most useful form of verification is indeed not confirming the third digit but confirming
that the gap we claim exists and has the right sign.

Third, parts 3 and 5 show that the protocol and the relations between events enter the
numbers. A host who opens doors at random lets only $66.472$ of
$100{,}000$ games run as far as the offer, and among those switching wins
$0.4998$ while staying wins $0.5002$; that gap of $0.0004$ is a quarter of a standard error.
The disjoint pair gives $P(A)P(B)-P(A\cap B)=0.1112$, matching $1/9$ to
three digits, while the independent pair gives $0.0001$, that is zero within random swing.
The number $0.1112$ is the size of the mistake made by anyone who thinks two disjoint events
must be independent.

What the program does not prove: not one of the five parts proves a formula.
The program only checks formulas that have already been proved, with a sample far larger
than the reach of the hand; Chapter 19 explains why that is legitimate.

\begin{untukdosen}
The teaching order I suggest: start from the $100{,}000$ table, not from
the axioms. Have the students fill in four cells: sick and positive, sick and negative,
healthy and positive, healthy and negative. The numbers $99$, $1$, $999$, and $98{,}901$
appear on their own, and from there $99/1.098$ becomes the conclusion of the class, not
the result of a lecture. After that the axioms, their consequences, counting, conditional
probability, Bayes, then independence. Monty Hall is best placed last, after the students
can write the host's protocol as a conditional event.

Two misconceptions appear most often. First, ``disjoint means independent''.
The symptom shows up in the arithmetic: students add the probabilities of two events that
should be multiplied, or the other way round, and then cannot explain why the result exceeds
one. Second, ``sensitivity equals the probability of being sick given a positive result''.
The symptom is the answer $0.99$ for the clinical problem, without using the prevalence
already given in the problem.

A blackboard activity that works: two students play the host with different written
protocols, one always opening a goat door and one opening at random; the class records
twenty games and then compares them with $2/3$ and $1/2$. Students are ready
to go further when they can give one example of two independent events and one example of
disjoint events with the same marginals, and compute a conditional probability from a
two-way table without reversing the direction.
\end{untukdosen}

\section{Summary}

\begin{ringkasan}
\begin{itemize}
  \item Probability is a function on events that is nonnegative, takes the value one on the
        whole space, and is countably additive on disjoint events; complement,
        monotonicity, and inclusion--exclusion follow from that.
  \item When outcomes are equally likely, probability is a ratio of counts, with the product
        rule and $\binom{n}{k}$; the assumption of equal likelihood has to be checked.
  \item Conditional probability $P(A\mid B)=P(A\cap B)/P(B)$ recomputes inside $B$
        and satisfies the three axioms; the product rule and the law of total probability
        are its consequences.
  \item Bayes' rule is not a new assumption but a definition read in two directions.
        Prevalence determines the answer more strongly than the quality of the test: with
        sensitivity and specificity $0.99$, a prevalence of $0.1$ percent gives
        $P(S\mid+)=0.0902$.
  \item Disjoint and independent are different; when both probabilities are positive,
        disjoint events cannot be independent, and only the intersection determines the
        relation.
  \item The answer to a paradox depends on the protocol and the unit of counting: Monty Hall
        $2/3$ under a fixed protocol and $1/2$ under a random one; birthdays $0.5073$ for
        $23$ people because what is counted is $253$ pairs. A seeded simulation
        checks, it does not prove.
\end{itemize}
\end{ringkasan}

\section{Exercises}

\begin{latihan}
\begin{enumerate}
  \item Prove the inclusion--exclusion formula for three events by applying the two-event
        version to $A\cup B$ and $C$. Show at which step the term
        $P(A\cap B\cap C)$ appears, and why its sign is positive.
  \item Assess this claim: ``Because $P(A\mid B)=0$, $A$ and $B$ are disjoint.'' True,
        false, or not enough evidence? Give an example with $P(A\mid B)=0$ but
        $A\cap B\neq\varnothing$, or explain why that is impossible.
  \item A test has prevalence $2$ percent, sensitivity $0.95$, specificity $0.90$.
        Compute $P(S\mid+)$ from a table of $10.000$ people, then repeat with
        specificity $0.99$; state both as fractions.
  \item Prove that if $A$ and $B$ are independent, then $A$ and $B^{c}$ are also
        independent. Use $P(A)=P(A\cap B)+P(A\cap B^{c})$ and explain why
        the same result holds for $A^{c}$ and $B^{c}$.
  \item Compute the probability that a pair with the same birthday exists for $30$ people,
        compare it with $1-e^{-\binom{30}{2}/365}$, then explain the source of the gap.
  \item Prove the law of total probability for $k=2$ directly from the axioms, without using
        the law for general $k$: write $A=(A\cap B)\cup(A\cap B^{c})$ and check
        that the two parts are disjoint.
  \item Four doors, one car, three goats. A host who knows the contents opens two
        goat doors among the three other doors and offers a switch to the one door
        that remains. Compute the probability of winning by switching, then check it with a
        simulation of $10^{4}$ games.
  \item Assess this claim: ``Two independent tests with sensitivity and specificity $0.99$
        give certainty $0.9999$ when both are positive, as long as the prevalence is $0.1$
        percent.'' Compute the correct number and state the assumptions.
  \item Is pairwise independence for three events enough to conclude
        $P(A\cap B\cap C)=P(A)P(B)P(C)$? Refute it with an example from two coin tosses,
        and compute both values.
\end{enumerate}

Hints for the three hardest items. Item 2: in a finite space where every outcome has
positive probability, $P(A\mid B)=0$ forces $A\cap B=\varnothing$, so for a die throw
the claim is true; a counterexample has to come from a continuous space, for instance
$A=\{1/2\}$ and $B=(0,1)$ with uniform probability, where $P(A)=0$ even though $A$ contains
a point. Item 7: at first the car is behind the chosen door with probability $1/4$; if not,
 the only door the host does not open must hold the car, so the probability of switching is
$3/4$. Item 9: take $A$ heads on the first coin, $B$ heads on the second coin, and $C$ both
coins the same; all three pairs give $1/4$ while the intersection is $1/4$, not $1/8$.
\end{latihan}

\begin{untukpeneliti}
What is genuinely proved in this chapter: the three axioms together with their consequences;
Bayes' rule as a theorem derived from the definition of conditional probability; the
impossibility of two disjoint events with positive probability also being independent; and
the dependence of the Monty Hall answer on the protocol, an analysis stated sharply by
Morgan, Chaganty, Dahiya, and Doviak, \emph{The American Statistician} 1991, in the article
``Let's Make a Deal: The Player's Dilemma''. What is popular but not proved here
is the claim that the natural frequency format improves prevalence reasoning in humans;
that is an empirical claim about readers, not a theorem, and its evidence lies in the
psychology of decision making literature.
\end{untukpeneliti}

Primary sources: Blitzstein and Hwang, \emph{Introduction to Probability}, Chapter 1 (Why
probability, Naive definition of probability, How to count, Story proofs, Non-naive
definition of probability) and Chapter 2 (The importance of thinking conditionally,
Definition and intuition, Bayes' rule and the law of total probability, Conditional
probabilities are probabilities, Independence of events, Coherency of Bayes' rule);
as well as Wasserman, \emph{All of Statistics}, Chapter 1 in the section Probability for the
formulation of the axioms and Chapter 2 in the section Independent Events.

Limits of validity: the whole chapter works in finite spaces where every outcome has
positive probability, where $P(A)=0$ is equivalent to $A=\varnothing$. In uncountable
spaces, conditioning on an event of probability zero is not defined by the ratio formula
and requires a regular conditional distribution; the Borel paradox shows that the answer
can depend on how the limit is taken. Independence is not stable either: it can disappear
or appear under conditioning, and in graphical models it appears as selection bias.

An open question: how far must the prevalence be known before Bayes' rule is used
honestly, and what can be said when the prevalence is known only from a small sample. The
formal answer is Bayesian inference with a prior, which brings in the problem of choosing
the prior; the frequentist answer is the hypothesis test of Chapter 20, which does not use
a prior but answers a different question. This chapter does not choose either.

\peta{Chapter 17 in Blitzstein and Hwang, \emph{Introduction to Probability}: Chapter 1
(Probability and counting, including How to count and the birthday problem example) for the
axioms, counting, and combinatorial probability; Chapter 2 (Conditional probability, Bayes'
rule and the law of total probability, Independence of events) for conditional probability,
Bayes, and independence. In Wasserman, \emph{All of Statistics}: Chapter 1 in the section
Probability and Chapter 2 in the section Independent Events. Sets and complements are used
from Chapter 2 of this book; the series underlying countable additivity from Chapter 5; the
law of large numbers that supports the simulation is proved in Chapter 18. Random variables
in Chapter 17 move the weight from events to numbers, while the hypothesis tests in Chapter
19 use Bayes' rule and the law of total probability for a different question.}
